\documentclass[a4paper,12pt]{article}

\usepackage[utf8]{inputenc}
\usepackage[T1]{fontenc}
\usepackage[italian]{babel}
\usepackage{amsmath}
\usepackage{amssymb}
\usepackage{graphicx}
\usepackage{hyperref}
\usepackage{geometry}
\usepackage{enumitem}
\usepackage{longtable}
\geometry{margin=1in}

\title{HiHome}
\author{Denise Nanni\\\href{mailto:denise.nanni@studio.unibo.it}{denise.nanni@studio.unibo.it} \and Juri Guglielmi\\\href{mailto:juri.guglielmi@studio.unibo.it}{juri.guglielmi@example.com} \and Nicolò Monaldini\\\href{mailto:nicolo.monaldini@example.com}{nicolo.monaldini@example.com}}
\date{Giugno 2026}

\begin{document}

\maketitle

\tableofcontents
\newpage

\section{Introduzione}
Il progetto HiHome consiste in un applicativo per gestire dispositivi all'interno di una smart home. Permette di gestire dispositivi smart da applicazione e aggiungere automazioni. Sono presenti funzionalità come notifiche, log di event relativi ai device smart, informazioni sullo stato attuale di temperatura, meteo e qualità dell'aria nella zona e grafici relativi ai valori passati di questi ultimi. Inoltre, è integrato con un agente AI in modo da effettuare alcune azioni in linguaggio naturale. L'applicazione supporta due tipi di utenti:

\begin{itemize}
    \item \textbf{Admin}: ha accesso a funzioni avanzate sulla gestione della casa, come log di eventi relativi a manipolazione di smart device, creazione di automazioni e aggiunta di device.
    \item \textbf{Utente Standard}: può agire sugli smart device, visualizzare statistiche relative allo stato attuale della casa, grafici relativi ad eventi passati, come lo stato della qualità dell'aria nella zona durante la settimana precedente e interagire con il chatbot.
\end{itemize}
\end{document}